\section{Introduction}
\label{sec:intro}

The motivating question for \sfs{} is simple: what should a filesystem look
like when many software agents edit one project concurrently across machine
boundaries?  Git provides durable history, branching, and explicit merge, but
its repository and commit model sits above the live filesystem.  File
synchronizers propagate live state, but commonly reduce a concurrent edit to a
whole-file conflict.  An \emph{agentic-time filesystem} should make history and
causality intrinsic, synchronize continuously, preserve rather than overwrite
unresolved alternatives, and still present ordinary files to tools that know
nothing about the replication protocol.

The immediate design prompt came from Theo Browne's 2026 discussion of source
control, filesystems, and a ``Dropbox for devs''~\citep{Browne2026Ideas}.  The
segment asks for one code-tree structure across local machines and cloud
agents, with content materialized when a participant touches it rather than
through manual clone and worktree choreography.  \sfs{} takes that prompt as a
filesystem problem, then adds the properties needed for concurrent autonomous
writers: retained causal history, fragment-level conflict visibility, client-
side encryption, and signed writer authorization.  We use \emph{agentic time}
for this resulting workload, not as a claim that the later mechanisms appeared
in the original prompt.

Local storage abstractions are usually divided by execution and trust boundary.
A kernel filesystem owns raw media and POSIX semantics; a portable userspace
filesystem places a container on another filesystem; a browser uses an
in-memory or origin-private backend; and cloud synchronization converts local
state into provider objects.  Snapshots or a version-control system add
history, while an encrypting layer protects the remote copy.  The resulting
stack works, but the boundaries duplicate identity, causality, indexing, and
recovery state.

This mismatch matters for local-first workloads~\citep{Kleppmann2019LocalFirst}
in which the same state is edited on several machines, may be stored by a
provider that should not see its contents, and is consumed through more than
one execution environment.  Agentic development is the primary workload:
small, frequent, partly overlapping changes from independent actors must remain
recoverable and synchronizable without forcing each agent to coordinate an
explicit commit.  The architecture also applies to financial, collaborative,
or personal-data applications whose virtual filesystem is useful only when its
bytes remain client-side encrypted and writer-authorized from browser or native
client through persistence and synchronization.  The service should transport
and retain state without ever receiving the content or metadata key.

Prior systems cover large parts of this space.  Versioning filesystems retain
past states~\citep{Santry1999Elephant,Soules2003CVFS,Peterson2005Ext3cow}; Coda,
Ficus, and Ivy support optimistic replication~\citep{Kistler1992Coda,
Reiher1994Ficus,Muthitacharoen2002Ivy}; Plutus, SiRiUS, SUNDR, and SPORC address
sharing over untrusted storage under different threat
models~\citep{Kallahalla2003Plutus,Goh2003SiRiUS,Li2004SUNDR,Feldman2010SPORC};
and WNFS and Peergos combine encrypted, portable, collaborative
storage~\citep{WNFSspec,Peergos}.  Ceph demonstrates one distributed object
substrate behind file, block, and object interfaces~\citep{Weil2006Ceph}.
Consequently, our claim is not that any one capability, nor the general idea of
multiple surfaces over one substrate, is new.

\sfs{} explores a narrower systems question: \emph{what follows if a
position-addressed fragment version map is the common state representation
across execution and trust boundaries?}  A unit has a stable UUID, a sparse
version vector, and an ordered map from fragment position to the causal dot of
the selected bytes.  A write replaces only the affected positions.  The same
map then provides four functions:

\begin{enumerate}[leftmargin=1.6em,itemsep=2pt]
  \item it is the materialized current version and the index into retained
    fragment history;
  \item it is the manifest used to identify missing or changed synchronization
    payloads;
  \item together with the version vector, it classifies concurrent updates per
    fragment, automatically merging disjoint changes and preserving overlapping
    ones as strains; and
  \item its logical identity is independent of physical location and backend,
    allowing the same container semantics in a raw-device kernel VFS, a
    userspace host-file VFS, an in-memory WebAssembly adapter, and an encrypted
    remote transport.
\end{enumerate}

Several supporting choices make this representation usable as a filesystem and
application VFS.  Deterministic fixed-size fragments give constant-time
offset-to-fragment addressing within a geometry band.  A double-buffered
authenticated header is the atomic publication point, while an optional
write-ahead log and bounded undo records cover asynchronous and in-place paths.
Metadata is always protected by AES-GCM; content is selectable per fragment as
GCM, XTS, or plaintext.  The secure application profile selects GCM content and
an owner-signed writer set.  In that profile the service receives no content
key, path, metadata plaintext, or content plaintext.  Account, object,
version-vector, size, timing, and access-pattern metadata remain visible.
Per-record Ed25519 signatures and the writer set authorize multi-writer record
frontiers.  These mechanisms are described in \cref{sec:design}.

The design is implemented rather than simulated.  A
\texttt{\#![forbid(unsafe\_code)]} Rust core, a FUSE adapter, and an independent
in-kernel C driver read and write the current v12 format.  The same engine is
also exposed through a C ABI and experimental in-memory WebAssembly bindings.
Synchronization runs unchanged against a blind service or a live peer
implementing the same transport contract.  Controlled functional scenarios
follow convergence through two and three replicas, hosted TLS, and a live peer.
The performance study separately compares the raw-device kernel path and the
container-on-ext4 FUSE path only with matching stacks (\cref{sec:eval}).

\paragraph{Contributions.}
\begin{itemize}[leftmargin=1.4em,itemsep=2pt]
  \item \textbf{A fragment-causal storage representation} in which the ordered
    fragment map and sparse version vector jointly serve local history,
    synchronization, and sub-file conflict classification
    (\cref{subsec:mvcc,subsec:sync}).
  \item \textbf{Execution- and topology-independent storage}: one v12 logical
    format is implemented by a raw-device kernel VFS, a userspace host-file VFS,
    the Rust engine and C ABI, and experimental in-memory WebAssembly bindings;
    one synchronization protocol operates against either a blind multi-tenant
    service or another live container, including propagation of concurrent
    frontiers, writer sets, and epoch-tagged key grants
    (\cref{subsec:transport,sec:impl}).
  \item \textbf{An atomic-publication, recovery, and allocation design}
    (\cref{subsec:commit,subsec:alloc}): a double-buffered, sequence-wins header
    commit plus fragment-granular copy-on-write and bounded undo records over a
    flat three-region allocator.  An optional WAL serves the asynchronous write
    path.  File content avoids recursive path shadowing in a global CoW tree;
    the catalog tries retain their own CoW cost.
  \item \textbf{A controlled artifact evaluation}
    (\cref{sec:impl,sec:eval}): deterministic multi-replica scenarios for
    merge, retained conflict, resolution, hosted sync, and live-peer sync, plus
    a fault-aware Linux throughput study that reports both the competitive
    native path and the remaining encrypted FUSE-read deficit.
\end{itemize}

A document/KV projection is included as an annex
(\cref{sec:annex-nosql}) rather than a primary
contribution.  It is useful evidence that the unit abstraction is not tied to
POSIX paths, but the paper's subject is the backend and its VFS, SaaS, peer, and
WebAssembly boundaries.

\Cref{sec:background} defines the representation and positions the design;
\cref{sec:design} gives its mechanisms and invariants; \cref{sec:impl,sec:eval}
describe the artifact and evaluation; and \cref{sec:related,sec:discussion}
separate the contribution from related systems and state its limits.
